% ============================================================
% 第6章 模型选择与参数设置
% 【负责人：建模线（6.1 基线已预填实测配置；6.2/6.3 待补充）】
% 赛题要求“模型选择、参数设置”为说明文档必备内容。
% ============================================================
\section{模型选择与参数设置}

\subsection{基线模型：L2 正则逻辑回归（已完成）}

按 2026-09-18 会议共识，选定成熟度高、易复现的 L2 逻辑回归作为首要基线，
为后续复杂模型提供对照基准。实测配置与结果：

\begin{table}[H]
\centering
\caption{逻辑回归基线配置（2.0 版数据实测）}
\label{tab:baseline-cfg}
\small
\begin{tabular}{ll}
\toprule
项目 & 配置 \\
\midrule
样本 & 500 台车（正类 126），\texttt{main\_20\_40} 主口径 \\
特征 & 9 个基线特征组（见 4.1 节） \\
正则 & L2，$C=0.1$（强正则抑制小样本过拟合） \\
求解 & Newton 法（各折 6 次迭代收敛，梯度 $\sim 10^{-12}$） \\
划分 & 车辆级分层 5 折（见 5.2 节），种子 20260917 \\
\bottomrule
\end{tabular}
\end{table}

\subsection{基线扩展与候选模型}

\begin{TODO}
负责人：建模线。按会议共识补充：(1) 随机森林基线（与逻辑回归并列的对照基准）；
(2) 候选模型阶梯（如 GBDT 类、以及是否引入 TabPFN 等外部预训练模型——注意
需先确认官方合规性，见附录 C 问题 7）；(3) 概率校准方法。全部实验须满足第 5 章
准入规则并登记配置。
\end{TODO}

\subsection{建模方案：单一模型与集成学习双线并行}

\begin{TODO}
负责人：建模线。按会议共识撰写双方案对比：
\textbf{方案 A（单一模型）}——全部数据混合训练一个通用模型，由模型自身学习
车辆类型差异（里程长短、有无盲区雷达等分布不均问题）；
\textbf{方案 B（集成学习）}——先按关键特征（里程区间、有无雷达等）聚类分群，
训练多个专用子模型，再通过动态路由或加权融合集成。
两方案并行开发，最终以验证集折外 AUC 与 Recall@K 择优（准入规则见 5.4 节）。
本节需给出：分群特征选择依据、融合方式、与方案 A 的对照实验设计。
\end{TODO}
